\section{From Blind Signatures to ARC}
\label{sec:arc-construction}

This section adapts the blind signature of Section~\ref{sec:blind-signature} into an anonymous rate-limited credential
(ARC) system.   To do so, $\Holder$ changes the input to $h$ as \(\vecmu\Vert \veck\), where \(\vecmu\) is the message and \(\veck\) is a uniformly sampled secret PRF key. It's easy to restructure our construction so that $\vecmu$ is instead a vector of private attributes, a function of which will be revealed by $\Holder$ during showing.

\paragraph{Rate-limiting via the secret key \(\veck\).}
Rate limiting is implemented by deterministic tags derived from a PRF keyed by the credential’s secret \(\veck\). For a
 nonce \(\nonce\in\mathcal{D}\), the holder computes
\[
  \prftag \;=\; \PRF(\veck,\nonce)
\]
and attaches it as part of the presentation. It additionally proves correctness of $\prftag$ with respect to $\veck$ and that $\nonce \in \nonceSpace$ i.e. in the space of permitted nonces. Intuitively, if the holder never repeats a nonce, then \(\prftag\) is pseudorandom and yields unlinkable showings.  If the holder repeats a nonce, the tag repeats deterministically, enabling the verifier to detect  double-spending for that nonce.
The verifier enforces its rate policy by maintaining state that records previously accepted \(\prftag\) values. 

\subsubsection{ARC protocol}
\label{subsec:arc-protocol}
We assume the same setup as in Section \ref{subsec:bs-protocol}. Moreover, define the relations 
\[
\relIssueARC:=\left\{ \left(  \begin{pmatrix}
	\ComPublicParam,\vecc,\\ 
    t, \vecr_0^{\Issuer},\vecr_1^{\Issuer}, \\
    \mA,B
\end{pmatrix},(\vecw=\begin{bmatrix}
\vecmu \\
\veck \\
\vecr_0^{\Holder}\\
 \vecr_1^{\Holder} \\
 \bar r
\end{bmatrix})\right) ~\bigg |~ \begin{matrix} \ComPublicParam=(\mACom,\BoundMsg)  \land  \\
\vecc = \mACom \vecw  \land \\
 ||\vecw||_{\infty} \leq \BoundMsg \land \\ 
\vecr_0= \centerfunc(\vecr_0^{\Holder}+\vecr_0^{\Issuer}) \land  \\
\vecr_1=\centerfunc(\vecr_1^{\Holder}+\vecr_1^{\Issuer}) \land \\
\vect = h_{\vecmu \| \veck,(\vecr_0,\vecr_1)}(B) \end{matrix}  \right\}
\]
and
\[
\relShow :=\left\{ \left( \begin{pmatrix}
\vecmu,\prftag,  \mA, \\
B, \BoundMsg,\BoundSig
\end{pmatrix},(\vecw = \begin{bmatrix}
\veck \\ \vecr_0\\ \vecr_1 \\ \vecu \\
\nonce
\end{bmatrix})\right) ~\bigg |~ \begin{matrix} 
  ||\vecmu,\veck,\vecr_0,\vecr_1||_{\infty} \leq \BoundMsg \land\\
  ||\vecu||_\infty \leq \BoundSig  \land \\
  \prftag = \PRF(\veck,\nonce) \land \\
  \nonce \in \nonceSpace \land \\
	\mA \vecu = h_{\vecmu\| \veck ,(\vecr_0,\vecr_1)}(B) \end{matrix}  \right\}
\]

Our Anonymous Rate-Limiting Credential scheme is then defined as follows:
	\begin{description}
		  \item[$\Setup(\SecParam)$ and $\IKeyGen$:] Same as Section \ref{sec:blind-signature}.
         
		\item[$\Issue$:] $\Issuer$ holds $\sk$ while $\Holder$ has $\mu,\vk$. 
        \begin{enumerate}
\item $\Holder$ samples \(\veck,\vecr_0^{\Holder},\vecr_0^{\Holder},\bar \vecr\) coefficient-wise in \([-\BoundMsg,\BoundMsg]\).
\item $\Holder$ computes \(\vecc\gets \mACom\cdot [\vecmu \Vert \veck \Vert  \vecr_0^{\Holder}\Vert \vecr_0^{\Holder}\Vert \bar \vecr]^\top\) and sends \(\vecc\) to \(\Issuer\).
\item $\Issuer$ samples \(\vecr_0^{\Issuer},\vecr_1^{\Issuer}\) coefficient-wise in \([-\BoundMsg,\BoundMsg]\) and sends \((\vecr_0^{\Issuer},\vecr_1^{\Issuer})\) to \(\Holder\).
\item $\Holder$ sets \(\vecr_0 \gets \centerfunc(\vecr_0^{\Holder}+\vecr_0^{\Issuer})\) and \(\vecr_1 \gets \centerfunc(\vecr_0^{\Holder}+\vecr_1^{\Issuer})\).
\item $\Holder$ computes \(\vect\gets h_{\vecmu,(\vecr_0,\vecr_1)}(B)\) and produces the
\emph{issuing} NIZK \(\nizkIssue\) using $\nizkscheme.\ZKProve$ proving knowledge of \((\vecmu,\veck,\vecr_0^{\Holder},\vecr_0^{\Holder},\bar \vecr)\) satisfying relation $\relIssueARC$.
\item $\Holder$ sends \((\vect,\nizkIssue)\) to \(\Issuer\).
\item $\Issuer$ verifies \(\nizkIssue\) for relation $\relIssueARC$ on $\vecc,\vect$ using $\nizkscheme.\ZKVerify$; if it accepts, sample \(\vecu \gets \TDRH.\SampPre(\trapdoor,\vect)\) and send \(\vecu\) to \(\Holder\). Otherwise abort.
\item $\Holder$ checks \(\mA \vecu = \vect\) and that \( || \vecu ||_\infty \leq \BoundSig  \). If yes, then output $\credential=(\vecu,\vecmu,\veck,\vecr_0,\vecr_1)$. If not, then abort. 

        \end{enumerate}
        \item[$\Show(\credential=(\vecu,\vecmu,\veck,\vecr_0,\vecr_1),\nonce)$:]~
        \begin{enumerate}
        \item Compute  $\prftag \;=\; \PRF(\veck,\nonce)$.
            \item Compute a 
\emph{showing} NIZK \(\nizkShow\) using $\nizkscheme.\ZKProve$ proving knowledge of \((\veck,\vecr_0,\vecr_1,\vecu,\nonce)\) for the relation $\relShow$.
\item Output $\presentation:=(\nizkShow,\prftag)$.
        \end{enumerate}
		\item[$\Verify(\vk,\presentation=(\nizkShow,\prftag),\mu,\algstate)$:] For an initially empty $\algstate$:
        \begin{enumerate}
            \item Verify $\nizkShow,\mu,\prftag$ for the relation $\relShow$ using $\nizkscheme.\ZKVerify$ and output $(0,\algstate)$ if the proof rejects.
            \item If $(\mu,\prftag)\in \algstate$ then output reject, otherwise add $(\mu,\prftag)$ to $\algstate$ and output $(1,\algstate)$.
        \end{enumerate}

	\end{description}

\subsection{Concrete PRF instantiation (Poseidon2-like)}
\label{subsec:poseidon-prf}

\providecommand{\Tr}{\mathrm{Tr}}

We instantiate the PRF from Definition~\ref{def:prf} using the Poseidon2-like construction of
Definitions~\ref{def:poseidon-prf} and \ref{def:poseidon-perm}. Algorithm~\ref{alg:prf} summarizes the computation.

\begin{algorithm}[H]
\caption{$\PRF(k,n)$ (Poseidon2-like with feed-forward + truncation)}
\label{alg:prf}
\begin{algorithmic}[1]
\\state $x \gets k\Vert n$; $x^{(0)}\gets x$
\For{$i=1$ to $R_F/2$} \state $x\gets E_i(x)$ \EndFor
\For{$i=1$ to $R_P$} \\state $x\gets I_i(x)$ \EndFor
\For{$i=R_F/2+1$ to $R_F$} \\state $x\gets E_i(x)$ \EndFor
\\state $y \gets x + x^{(0)}$ \Comment{feed-forward}
\\state \Return $\Tr(y)$ \Comment{first $\ell_{\prftag}$ lanes}
\end{algorithmic}
\end{algorithm}

\paragraph{Proof-friendly arithmetization.}
In the post-sign NIZK we enforce $\prftag=\PRF(\veck,\nonce)$ using a checkpointed arithmetization: we commit
selected S-box outputs and let the verifier replay all linear wiring between checkpoints from openings and public
constants. Section~\ref{subsec:prf-checkpointed} states the constraint families; Appendix~\ref{app:prf} records the
checkpoint schedule and extended notes.

\begin{comment}
Preserved from Carsten's 10-04 source: this entire PRF-instantiation subsection was wrapped in a comment environment at
that revision boundary. The active paper keeps the subsection live because the later rewrite uses it to anchor the
SmallWood PRF constraint discussion.
\end{comment}


\begin{comment}
\subsubsection{ARC protocol}
\label{subsec:arc-protocol}
We assume the same setup as in Section \ref{subsec:bs-protocol}. Moreover, define the relations
\[
\relIssueARC,\qquad \relShow.
\]
In the 10-04 source, \(\relIssueARC\) explicitly committed \((\vecmu,\veck,\vecr_0^{\Holder},\vecr_1^{\Holder},\bar r)\)
and required
\[
\vect = h_{\vecmu \Vert \veck,(\vecr_0,\vecr_1)}(B),
\]
while \(\relShow\) explicitly bound
\[
\prftag = \PRF(\veck,\nonce),\qquad
\mA \vecu = h_{\vecmu\Vert \veck ,(\vecr_0,\vecr_1)}(B).
\]
\cb{FIX THIS}
That version then presented explicit \(\Issue\), \(\Show\), and \(\Verify\) routines in terms of
\(\relIssueARC,\relShow\), prior to the later theorem-clean rewrite that compiles showing through the Section~5
SmallWood relation instead.

\subsection{Showing protocol}
The 10-04 source also stated an older abstract showing algorithm in which \(\pi_{\mathsf{post}}\) directly proved
knowledge of \((u,\mu,\veck,r_0,r_1)\) satisfying the signature/hash relation, the PRF-tag relation, and the public
bounds, rather than referring forward to the compiled Section~5 relation.
\end{comment}



\subsection{ARC correctness and conditional security decomposition}

\paragraph{Honest correctness.}
If issuance completes, the holder stores a witness \((u,\vecm,\veck,r_0,r_1)\), together with
\(\vecmu=\vecm\Vert \veck\), satisfying the abstract signature relation
of Section~\ref{sec:blind-signature}. For any public nonce \(\nonce\), the holder computes
\(\prftag = \PRF(\veck,\nonce)\) and derives the corresponding witness for the compiled post-sign relation of
Section~\ref{sec:smallwood-model}. By completeness of the post-sign NIZK, the verifier accepts any honest presentation
for which the local policy check on \((\nonce,\prftag)\) also accepts.

\paragraph{Cryptographic decomposition of the ARC layer.}
The ARC layer combines five ingredients:
\begin{enumerate}[leftmargin=1.25em]
    \item the non-blind theorem-backed signature layer of Lemma~\ref{lem:vsis-unforgeability};
    \item blindness and one-more unforgeability of the blind issuance protocol, which remain to be reduced;
    \item straightline-extractable knowledge soundness and zero knowledge of the post-sign NIZK for the compiled relation of Section~\ref{sec:smallwood-model};
    \item pseudorandomness and correctness of the PRF; and
    \item verifier-side policy state recording previously accepted \((\nonce,\prftag)\) pairs.
\end{enumerate}
Only the first and third ingredients are theorem-backed in the present draft. The fourth is an external assumption on
the chosen PRF instantiation, the second remains a reduction goal, and the fifth is an application-layer enforcement
rule rather than a stand-alone cryptographic theorem.

\paragraph{Conditional policy soundness.}
Knowledge soundness of the post-sign proof yields a witness for the compiled showing relation of
Section~\ref{sec:smallwood-model}. Section~\ref{subsec:statement-strength} is the sole bridge from that compiled
replay-space statement to the coefficient-domain statement used here. In particular,
Corollary~\ref{cor:showing-proof-consequence} yields a witness \((u,T,M,K,R_0,R_1,Z)\) satisfying
\[
B_3Au-(Au)R_1-B_0-B_1(M+K)-B_2R_0=0
\]
in \(R_q\), together with
\[
(B_3-R_1)\Had Z = 1,
\qquad
Au = h_{M+K,(R_0,R_1)}(B).
\]
Under that signature-side condition, together with one-more unforgeability of the blind issuance layer and correctness
of the PRF, the signed key \(\veck\) determines the public tag uniquely for each nonce. Consequently, a verifier that
records spent \((\nonce,\prftag)\) pairs rejects same-nonce replays deterministically.

\paragraph{Conditional presentation unlinkability.}
Presentation unlinkability in Definition~\ref{def:pres-unlink} is not implied by PRF security and post-sign zero
knowledge alone. In that experiment the adversary controls the issuer during the two issuance sessions. Accordingly,
unlinkability also depends on an issuance-layer privacy property, namely blindness or a comparably strong issuer-view
hiding statement. Under that issuance-layer privacy assumption, post-sign zero knowledge hides the witness used during
showing, and PRF pseudorandomness makes the public tags computationally random across distinct fresh nonces. This yields
the intended unlinkability claim for showing transcripts, subject to the explicit nonce-reuse caveat.

\paragraph{Attribute privacy and selective disclosure.}
The message part \(\vecm\) can encode attributes; the post-sign NIZK can be extended to prove predicates about \(\vecm\) and
to selectively reveal components, without changing the issuance protocol. We focus on the core rate-limiting
construction and defer predicate engineering to the SmallWood programming model in Section~\ref{sec:smallwood-model}.

\begin{comment}
\subsection{ARC security sketch}
\paragraph{Correctness.}
If issuance completes and the holder uses a fresh nonce, then the post-sign statement is satisfiable by construction.

\paragraph{Unforgeability and policy soundness.}
By knowledge soundness, any accepting presentation yields a witness satisfying both the signature/hash relation and the
PRF-tag relation. Under one-more unforgeability of the blind signature, an adversary cannot produce accepting showings
for credentials it did not obtain via issuance. Policy soundness is enforced by verifier state because
\(\prftag=\PRF(\veck,\nonce)\) is uniquely determined by \((\veck,\nonce)\).

\paragraph{Unlinkability.}
Presentation unlinkability is formalized in Definition~\ref{def:pres-unlink}. For distinct fresh nonces, PRF security
and zero knowledge hide the credential secret, except for the intended nonce-reuse caveat.

\paragraph{Constraint view (what the post-sign NIZK must prove).}
To prove \(\mathsf{tag}=\PRF(k,\mathsf{nonce})\) inside SmallWood, where \(k\) is the signed secret part of the
message, we arithmetize the S-box constraints, the public linear layers, and the feed-forward plus truncation
constraints binding the public tag to the final state.
\end{comment}
